% Further findings for n=5; scope is deliberately explicit.
\section{K-010: recursive spectral certificates and the first five-dimensional classification}

For three multisets $W=(A,B,C)$ of $n$ residues modulo a power of two $M$,
write $d(W)=\sum A+\sum B+\sum C\pmod M$.  An aggregate edge is a triple
$(a,b,c)$ of distinct residue values, one from each multiset, with
$a+b+c=0\pmod M$.  Its three margin targets are the multiplicities of the
corresponding residue values.  We call $W$ \emph{torus feasible} when these
aggregate edges admit nonnegative rational masses with those margins.
By the argument in K-002, this is a necessary condition for the fixed
space to contain a nonzero semistable tensor.  Repeated eigenvalues are
included, and torus feasibility is not used as an unqualified substitute
for semistability.

\subsection*{The refinement lemma}
Every torus-feasible triple modulo $2M$ reduces to a torus-feasible triple
modulo $M$.  Indeed, sum its masses over the residue classes that merge
under reduction; both the edge equation and each prescribed margin
reduce to the required equations.  Conversely, all possible lifts of a
fixed residue multiset are obtained as follows.  If a residue $a$ has
multiplicity $m_a$, put $j$ of its copies at $a+M$ and the other $m_a-j$
copies at $a$, independently for $0\le j\le m_a$.  Thus a factor has
$\prod_a(m_a+1)$ possible refinements.

We may quotient at each stage by coordinate permutations, factor
permutations, scalar shifts whose three residues sum to zero, and odd
units on all exponents, with the latter chosen to normalize the prescribed
determinant exponent.  This quotient does not lose fine-level cases.
Coordinate and factor permutations lift immediately.  An odd unit modulo
$M$ has an odd lift modulo $2M$.  Scalar shifts with sum zero modulo $M$
have lifts with sum zero modulo $2M$, by adding $M$ to one shift if
necessary.  Applying these lifted operations to a fine-level triple gives
a refinement of its normalized coarse representative.  The independent
choices just described include every such refinement.

For a target exponent $q$ modulo $M$, the canonical representative is the
lexicographically least triple obtained by all odd units $u$ satisfying
$u d(W)=q\pmod M$, all supported-edge shifts, and factor permutations.
A supported-edge shift subtracts the three coordinates of an aggregate
edge, making zero occur in every factor.  Every feasible support has an
edge, so these shifts suffice.  If the target is primitive, the normalizing
unit is unique.  If it is nonprimitive or zero, \emph{all} normalizing
units are retained.  The base case modulo two lists all multisets
containing zero and all unordered triples of these multisets.  These
observations prove completeness of the recursion, without assumptions
on eigenvalue multiplicities or on faithfulness of the determinant
character.

\subsection*{Exact certificates for primitive characters}
The Python implementation \texttt{recursive\_spectra.py} applies this
recursion to primitive targets.  Every rejected triple comes with a
rational Farkas certificate: for its aggregate incidence matrix $A$ and
margin vector $b$, the recorded vector $y$ satisfies
\[
                       A^t y\ge0,\qquad b^t y<0.
\]
Every accepted triple comes with a rational vector $x\ge0$ satisfying
$Ax=b$.  The independent checker
\texttt{verify\_recursive\_certificates.py} uses only Python integers and
\texttt{Fraction}; it checks these equalities and inequalities directly
and checks the resulting counts against the recorded classification.
For $n=5$ the complete counts are
\[
\begin{array}{c|r|r|r|r}
 M&\text{raw refinements}&\text{canonical candidates}&
    \text{feasible}&\text{infeasible}\\\hline
 2&35&7&3&4\\
 4&4168&154&30&124\\
 8&210456&10447&375&10072\\
 16&8394752&243824&77&243747\\
 32&2514944&36878&0&36878
\end{array}                                                    \tag{K-010a}
\]
All $291310$ rational certificates pass that independent checker.
There are therefore exactly $77$ torus-feasible normalized spectral
types of order-sixteen lifts with primitive determinant character,
and no primitive-character lift of order $32$ or any larger power of
two.  The assertion about larger powers follows by reduction to modulus
$32$; it does not assume that all determinant characters have faithful
lifts.  Combining this with K-009, whose five-dimensional two-primary
lift cutoff is $64$, excludes determinant character order $64$.
At this stage alone determinant order $32$ can still arise from an
order-$64$ lift, and determinant order $16$ can still arise from
order-$32$ or order-$64$ lifts.  The subsequent prescribed-character
search addresses exactly those possibilities.

The complete list of the $77$ types, their positive masses and their
normalized spectra is in
\texttt{recursive\_n5\_m16\_result.json}; the accompanying
\texttt{N5\_SPECTRA\_TABLE.tex} prints all the spectra.  This is a
finite-lift classification, not yet a count of distinct quotient curves.

\subsection*{Polystability of all seventy-seven types}
Seventy-six types have simple spectra in all three factors.  Each has
exactly thirteen fixed tensor coordinates, a rank-thirteen incidence
matrix, and a strictly positive, unique margin-one vector.  The exact
Smith diagonals are $(1^{12},16)$.  These statements are checked by
\texttt{analyze\_n5\_types.py} and recorded in
\texttt{n5\_type\_analysis.json}.  With coefficients equal to the square
roots of the positive masses, the three reduced matrices are identities.
Off-diagonal entries vanish because distinct fixed support edges cannot
agree in two coordinates when all three spectra are simple.  Thus the
tensor is critical and its $SL_5^3$ orbit is closed.  The coefficient-one
tensor lies in the same product-linear orbit: the incidence map has full
column rank, so the complex diagonal torus acts transitively on the
nonzero coefficients on this support.  Necessity of the unique positive
mass vector also shows that every semistable tensor in this fixed space
has all thirteen coefficients nonzero.  Consequently there is exactly
one product-linear orbit of semistable tensors for each fixed spectral
triple, before possible coincidences between different spectral triples.

The remaining type, the first in the list, is
\[
 A=(0,0,1,6,13),\quad B=(0,2,3,7,13),\quad C=(0,3,8,12,13).
                                                               \tag{K-010b}
\]
Its aggregate support, in lexicographic order, is
\[
\begin{gathered}
 (0,0,0),(0,3,13),(0,13,3),(1,2,13),(1,3,12),(1,7,8),\\
 (6,2,8),(6,7,3),(6,13,13),(13,0,3),(13,3,0),(13,7,12),
\end{gathered}
\]
and its unique positive aggregate mass vector is
\[
 \frac1{16}(13,7,12,5,6,5,11,1,4,3,3,10).                    \tag{K-010c}
\]
The first three entries correspond to three vectors in the repeated
two-dimensional eigenspace of the first factor.  Use the three columns
of
\[
 F=\begin{pmatrix}
 \sqrt3/2&-1/2&0\\[2pt]
 1/4&\sqrt3/4&-\sqrt3/2
 \end{pmatrix}.                                               \tag{K-010d}
\]
Their squared norms are $13/16,7/16,12/16$, and $FF^*=I_2$.
Give the other nine coordinates the positive square roots of their
listed masses.  No two aggregate edges agree in two residue coordinates,
and the three columns in (K-010d) occupy different second and third
coordinates.  Hence all off-block reduced-matrix entries vanish, while
the two-dimensional block is $FF^*$.  The three reduced matrices are
identities, proving polystability of this type as well.

For completeness, its semistable locus is also one orbit of the
centralizer of the prescribed diagonal element.  A zero aggregate
coordinate contradicts uniqueness and positivity of (K-010c).  If two
of the first three vectors were collinear, put their common line on
the first coordinate axis.  Their two aggregate masses would already
sum to more than one: the three pairwise sums are $20/16,25/16,19/16$.
No margin-one mass vector could then exist in that basis, contradicting
torus semistability.  Thus all three projective lines are distinct.
$PGL_2$ is transitive on ordered triples of distinct points of $\mathbb P^1$,
so these three vectors can be put on $e_0,e_1,e_0+e_1$.  The resulting
thirteen nonzero tensor coordinates have incidence rank thirteen;
diagonal scaling makes all their coefficients one.  This gives the
claimed orbit and also shows directly why repeated spectra cannot be
discarded.

For every one of these seventy-seven coefficient-one representatives,
the infinitesimal $GL_5^3$ orbit map has rank $73$.  The calculation
records a nonzero rank-$73$ minor modulo the prime $1009$.  Its rational
rank is therefore at least $73$, and the two-dimensional scalar kernel
makes it at most $73$.  Thus the connected full stabilizer is exactly
the scalar kernel.  This proves finite effective stabilizers, but it
does not prove that their finite determinant kernels are trivial.
The moving-normal determinant products recorded in the same file are
therefore explicitly labelled as products \emph{before a possible
finite quotient}; they are not asserted to be local Hilbert residues.
In particular no seventy-seven-term global residue sum is inferred
from this calculation alone.

\subsection*{Scheme-theoretic dictionary and remaining gate}
The incidence columns define a morphism of split tori over $\mathbb Z$.
Smith form $(1^{12},16)$ says that the finite part of its diagonal kernel
is the diagonalizable group scheme $\mu_{16}$, after the scalar torus is
separated.  Refinement from $M$ to $2M$ is a finite search on characters
of diagonalizable group schemes; the mass equations express the
Hilbert--Mumford convex-hull criterion over characteristic zero.
The critical representatives identify closed complex geometric orbits.
At this stage the further issue is the full finite group scheme
stabilizer and its action on the Luna slice, followed by the
multiple-component Koszul localization of P-F23g.  Infinitesimal
freeness alone does not settle that finite-group question.
K-012--K-013 resolve it for all seventy-seven types, K-011 supplies
the higher-lift coverage, and K-014 computes the resulting global pole.
